\documentclass[11pt]{article}
\usepackage[margin=1in]{geometry}
\usepackage{amsmath,amssymb,amsthm,booktabs}
\newtheorem{theorem}{Theorem}
\newtheorem{proposition}{Proposition}
\title{Relative Tree Gauge and Boundary-State Disintegration}
\author{}
\date{}
\begin{document}
\maketitle

\begin{abstract}
We give an exact finite-graph separation of an $SU(2)$ lattice connection into
a boundary-covariant retained state and an interior conditional fiber.
Interior gauge redundancy is removed by a boundary-rooted forest, while the
boundary gauge action remains visible.  This supplies the typed interface
needed by state-weighted defect transfer.
\end{abstract}

\section{Relative tree gauge}
Let $\Gamma=(V,E)$ be a finite oriented graph with nonempty boundary
$\partial V$.  The relative gauge group is $G^{V\setminus\partial V}$, acting
on $U\in G^E$ by
\[
 U_{xy}\longmapsto g_xU_{xy}g_y^{-1},\qquad g_b=1
 \quad(b\in\partial V).
\]
Choose a forest $F$ containing one directed path from a boundary root to every
interior vertex.

\begin{theorem}[Relative tree-gauge coordinates]
For every $U\in G^E$, there is a unique relative gauge transformation that
sets every forest link to the identity.  It induces a bijection
\[
 G^E/G^{V\setminus\partial V}\cong G^{E\setminus F}.
\]
The normalized nonforest variables are invariant under interior gauge
transformations and covariant under independent transformations of their two
boundary roots.
\end{theorem}
\begin{proof}
Set the normalizer at every boundary root to one and propagate it recursively
along $F$: if $v$ is the child of $u$, put $h_v=h_uU_{uv}$, reversing the link
when required.  Then $h_uU_{uv}h_v^{-1}=1$.  A forest has a unique root path,
so the recursion and the normalizer are unique.  Under an interior gauge
$g$, the path product changes to $h_vg_v^{-1}$, which cancels in each
normalized nonforest edge.  If boundary gauges are allowed, the two root
factors remain on the left and right.
\end{proof}

\begin{proposition}[Haar disintegration]
With normalized Haar measure,
\[
 \prod_{e\in E}dU_e
 =\left(\prod_{v\in V\setminus\partial V}dg_v\right)
   \left(\prod_{e\in E\setminus F}dH_e\right).
\]
Consequently any local plaquette density, including the heat-kernel density,
has an exact conditional disintegration after the nonforest variables are
partitioned into retained boundary variables $\sigma$ and interior fiber
variables $\phi$.
\end{proposition}
\begin{proof}
The recursive change of variables uses only group multiplication and
inversion.  Haar invariance makes each triangular step unit-Jacobian.  Fubini's
theorem then gives the stated product disintegration.
\end{proof}

\section{Sufficiency}
For the reflection-symmetric $2\times2$ square, there are nine vertices,
twelve links, four plaquettes, eight boundary links, one forest link, and three
fiber links.  A forest invariant under horizontal reflection makes relative
normalization commute with that reflection.

The complete retained boundary connection determines every plaquette term
crossing from the block to its exterior.  Omitting a retained crossing edge is
not sufficient: changing only that edge changes the holonomy, and hence any
nonconstant central plaquette weight.  Center sign, cut-locus proximity, and
central versus regular stabilizer type are derived from retained Wilson loops;
they are not discarded labels.

\section{Scope}
The finite representation theorem is exact.  A complete $2^4$ block,
overlap-consistent coarse variables, a state weight, and a subcritical
nonlinear bad kernel remain open.  Reflection compatibility here does not by
itself prove reflection positivity of an iterated RG map.
\end{document}

